\chapter{---STRIKES.} \label{STRIKES}

When a body of men stand out for a price which their employers refuse to
give, while this dispute is pending, the position of the workmen is that
of a strike. As strikes are the last resort, as they are always
expensive, and as they engender mutual ill feeling, they should never be
entered into without duly calculating the probabilities of success, nor
until all means of amicably settling, the difference have failed. It
often happens that workmen have no alternative, but either to submit to
a reduction tyrannically enforced, without any reasoning on the matter
being allowed, or to cease from labour. Often has a strike thus been
precipitated, and ruin inflicted on employer and employed, which might
have been averted by a little calm reasoning on the matter. It is the
same when a rise in wages is asked by the men. Both parties are apt to
view each other as enemies, and in this jaundiced view, which prevails
equally on both sides---aggravated by the unconciliatory tone which is
sure to result from such a state of feeling---reasoning on the
subject---as the subject, considering the important results to both
paHies which are then pending, should be reasoned upon---is rendered
impossible. The beginning of strife is like the letting out of water
that might be, at the commencement, easily stopped. But if there be one
thing more than another which, in their turn, both parties in these
circumstances often, to all appearance, agree in throwing aside, it is
the conciliatory spirit which might prevent these ruinous disputes. But
while strikes are always to be deprecated, because they are, for the
time, a state of moral warfare, and, like all sates of hostility,
productive of mutual bitterness---and because they are carried on at a
loss to both parties---we are, notwithstanding, clearly of opinion, from
long experience of their results to journeymen both of success and
defeat, that there is no proper alternative, in certain cases, than the
position of a strike.

The following extract from a very talented article in the
\textit{Builder} of August 11, 1859, which very clearly sets forth the
nature of strikes, as clearly exhibits their rationale:---

\begin{quote}
    \hspace{2em} Suppose a question of cotton and sugar value instead of
    labour value---how does the seller know whether he is selling too
    cheap, except by refusing to sell at all below a certain higher
    rate? If there is but little cotton or sugar, as of course he
    suspects or affects to suspect there to be, he will sell what he
    chooses at his own figure; but if not, he must take the buyer's
    price for it. Now, the fixed price about which the transaction halts
    is the strike of the seller against the buyer---of the supplier
    against the demander---and provides the only practicable means of
    arriving at the fair value.
\end{quote}

Against strikes it has been often urged---

\begin{enumerate}
    \item Their great expense;
    \item That they promote the introduction of machinery, and
    consequently leave the workmen in a tenfold worse position than they
    were before.
\end{enumerate}

1. It is admitted that strikes are very expensive. But the expense of
anything must be taken in reference to the gain it is intended to
procure, or the loss it is intended to avert. It should also be viewed
in reference to its result in success or defeat.

The expense of strikes to prevent a reduction of wages, let it be what
it may, is soon equalled by the amount it is intended to take from the
men's wages. A reduction of a penny per hour---we take this sum, because
in our own trade reductions have been attempted which would amount to
it, but any other sum can be taken, as the case may be---in a day often
hours, occurring in a trade of 1,000 men, amounts to £250 a week, or
£13,000 a year.\footnote{At 5 per cent., £13,000 a year represents a
capital of £260,000} This may be called a small trade. In a trade of
10,000 men it amounts to £130,000 a year;\footnote{At 5 per cent.,
130,000 a year represents a capital of £2,600,000} and in a trade of
20,000 men it amounts to £260,000 a year.\footnote{At 5 per cent.,
£260,000 a year represents a capital of £5,200,000} Considering the
capital these sums per annum represent, it cannot be surprising that
workmen are willing to incur a very great expense to prevent loss so
enormous.

But suppose the men to be defeated after incurring great expense---the
employers' expenses are sure to be fully equal to those of their men,
besides the possible loss of business. The fact of its being often very
expensive to reduce wages, prevents reductions being attempted which
otherwise would be made without hesitation or scruple. Strikes,
therefore, even in defeat, have a powerful tendency to prevent a future
and further reduction of wages.

Besides, it might as well be urged because in a comparatively short
period millions have been expended in going to law, involving often
total ruin to families, loss of reason, and even suicide, that therefore
no unjust aggressive claim should be opposed, and referred to legal
decision.

With respect to the objection, that strikes promote the introduction of
machinery: at no time is a strike needed, or indeed is there any
stimulus, but the shortening of labour, necessary for its introduction,
whenever it is found to answer. It may be true that ``mules'' and
``double deckers'' were introduced during, or about the time of a
strike; but these, if they saved labour, would have been used if there
had been no strike. The spinning jenny was not invented in consequence
of a strike; neither was the power-loom. No strike existed when the
printing machine was introduced; nor was there any to occasion its
subsequent improvements. Nor has been, or would, the introduction of
machinery in any trade be delayed an hour beyond the time it was found
to answer. ``The ingenuity of our mechanics,'' alluded to in the
``Quarterly Review'' (p. 504) has never needed the impulse of a strike
to call it into action; indeed, the fact has been exactly the contrary.
It has been more apt to employ itself in the making of machines which
were found not to pay, than to require any stimulus, whether of a strike
or otherwise, to call it into action.

Two great mistakes are often committed on both sides in strikes. The
first is, in forgetting that the issue joined is simply to prove which
of the two parties can stand out longest in a bargain---that the dispute
is not contending as in a combat, but simply for the adjustment of a
certain price for labour. It is quite true, that, where a reduction of
wages results from one employer underselling other employers, and who is
seeking to make up his profit out of his men's wages, a sense of wrong
is felt by the men so reduced, similar to that which prompts a man to
repel an injury. And, doubtless, in such a case they think that such an
employer is seeking to extend his trade by taking what does not belong
to him, but to themselves. And, on the other hand, where a rise in wages
is asked, which the employer feels will reduce his profits, and much
narrow his operations, he also will feel the same. Nevertheless, to
bring such a contest to an issue at the least possible cost to both
parties, and the least injury to the trade itself, such feelings ought
to be laid aside, and the dispute conducted according to what it really
is, simply a bargain, in which the employer is the buyer, and the
employed is the seller. Both are standing out against each other's
price. If the seller in this case can do without selling, and the buyer
can get no labour elsewhere, the workman is sure to win. If the
contrary, the employer. It is indispensable, therefore, as the dispute
cannot be carried on without great loss to both, that each party should
be well informed of every circumstance connected with that their true
position, both at the outset and from day to day as the dispute
proceeds. This, however, is not always done, and is the other mistake
often committed. In the builders' strike and ``lock-out,'' if the
masters had at all studied their position as to the number of men likely
to be available in the precipitated strike at Messrs. Trollope's, they
would have found, if their own statements are to be depended on, that
there was no occasion whatever for their lock-out, with its great
expense to themselves, and suffering to their men. It was, as they
expressed it, ``to continue for a month, or until the Messrs. Trollope
and Sons were enabled to resume their operations.''

This firm, by its own statement, was enabled to resume operations in a
fortnight, without having gained over a single man that had been locked
out, either of their own late workmen or those of any other of the
combined employers; and it was the same at the subsequent period, when
they announced that they had obtained their full complement of men. And,
as the lock-out, therefore, did not achieve this, it was indisputably
not needed at all. Their loss of £130,000, as announced in the
\textit{Weekly Dispatch} of January 1, 1860, which, from circumstances
not necessary to allude to here, is unquestionably good authority, has
therefore been incurred without the smallest occasion for it. Many
blunders have no doubt been made by workmen in their strikes, but none
so great as this, because so entirely without real necessity.

Both parties, however, are too often blinded by passion, to see clearly
their way, and being angry, the contest is transformed from what it
really is, merely of a bargain, into that of a combat. Under the
influence of this excitement, the leaders of both parties are too often
more intent upon setting forth the injuries which their respective
audiences imagine they shall sustain, if they abate the least of their
original views; or, what is more mischievous, they feel it a duty to
urge that they cannot without \textit{disgrace} modify, for the purpose
of arrangement, the several points in dispute, forgetting all the while
that it is only a bargain for which they are contending, in which it is
the special business and duty of the seller, if he finds he cannot get
his price, and the market is being, or likely to be, supplied, to take
the next in value, and in no case to stand out, if he sees the market
being supplied without him. Among dealers, he who loses by standing out
too long is always an object of derision. This mistake is too often made
in a strike. The dogged valour which has so often enabled the sons of
toil to conquer in battle is not unfrequently the cause of total defeat
in a strike, when, with a little address, considerable advantage might
have been gained, or loss prevented. Such address has been considered
cowardly and disgraceful, and the affair has ended with the loss of
everything. On the other hand, the employers, under the influence of the
same feelings and mistake, have often been led to refuse all adjustment
or arbitration, and thereby have led, by protracting the struggle, to a
much greater loss than their victory has achieved gain; and in some
cases, as in the master-builders' lock-out, to immense loss, without any
substantial gain at all, unless the gratification of their anger be so
considered.

But, after all, the true position of employer and employed is that of
amity. They are each, notwithstanding these occasional disagreements,
the truest friends of the other, and neither can inflict an injury on
the other without its recoiling on himself. Capital and Labour should go
hand in hand. Experience has amply proved that the Capitalist cannot
injure the Labourer, or the Labourer the Capitalist, without each
inflicting injury, and perhaps ruin, upon themselves.

\section{THE ``EDINBURGH REVIEW.''}

Many of the charges by this reviewer have already been answered. It is
necessary, however, briefly to notice the article. It is written,
apparently, \textit{to order}---confessedly, in profound ignorance of
the subject. This ignorance may be real or feigned---feigned, in order
to induce a belief that Trades' Unions were secret societies, similar to
those which ``honeycomb continental nations,'' abounding in ``spies,''
with ``forced dumbness,'' and ``\textit{surveillance}'' and ``distrust
of neighbours on the right hand and left,'' of which no information
could be with certainty obtained. This, indeed, he largely insinuates,
which is turning his ignorance, if real, to very good account. And it
must be confessed, if this ignorance be feigned, that the acting, if it
be not altogether successful, is very effective. Knowing nothing, or
being obliged, from the course he has adopted, to pretend to know
nothing of the subject, and finding that Adam Smith and McCulloch were
dead against him, and Stuart Mill, who was formerly so, having only
recently modified his opinion and even now enjoining that there should
be no encouragement given to the multiplication of labourers, because in
the opinion of ``farmers,'' ``Boards of Guardians,'' and others of the
employer class, the checking of such multiplication would make them
``too independent,'' and therefore, besides in other respects, not a
presentable authority for his intended article, he with admirable
adroitness tells his readers at the commencement that he is ``not going
to preach political economy.''

Nor was the information respecting Trades' Unions that could be obtained
from ``Blue Books,'' namely, the Parliamentary Enquiry on Combinations of
1838, and the Report of the Hand-loom Weavers' Commission suitable for
his purpose; the last, though quoted by him at the end, not being
applicable to any extent, as the hand-loom weavers were not Trades'
Unionists, which is one reason given in the before-mentioned Enquiry of
1838 (Questions 1418--19) why their wages were so wretchedly low; while
the evidence in the same Enquiry (of 1838), which is divided into two
parts, the first containing chiefly the evidence on the cotton-spinners'
strike of 1837, and the second the evidence on Irish Trades' Unions, is
too conflicting to be used by him, who only intended to see the
employers' side of the question. Indeed, it is so conflicting that the
Committee made no report upon it. Besides, the first of these contained
the opinions of Sir Archibald Alison on Trades' Unions, whom no one will
suspect of having an undue leaning towards them. Apart from intimidation
and violence which we are as much against as he could possibly be he
expresses himself several times favourable to trades' combinations, in
terms similar to the following extract. He is speaking of the
CottonSpinners' Union of Glasgow:---

\begin{quote}
    \hspace{2em} ``Their combination, I think, might have a very
    beneficial effect in their own favour, if it was limited to merely
    legal acts. I think, for example, that it is a very good thing for
    the cotton-spinners, as well as every other class of labour, to
    combine, because it enables numbers, to a certain degree, to
    compensate and to enter with equality into the lists with capital;
    and therefore I think that combinations are essential to support the
    rights of labour in the competition with capital.'' (Question 1956.)
\end{quote}

In this state of things, where it was so essential to shut himself out
from using, if he possessed it, any knowledge of the subject, there was
nothing for it but to follow the course pursued by Lord Chesterfield
when he brought the Bill into the House of Lords for reforming the
Calendar. His lordship tells his son, in one of his Letters, that having
to bring in this Bill; knowing nothing of astronomy, and consequently
unable to explain the reasons for the change; he did what he believed
his audience would be better pleased with, namely, instead of urging the
reason for the alteration from the science---which, he further said, he
did not believe his audience would have understood if he had been able
to give such explanation---he gave a history of the changes made in the
Calendar down to the then present time, taking care to flatter the
prejudices of his hearers by telling them what he thought they would
like to hear, by which they were so delighted that they gave him credit
for his profound knowledge of the science; while another noble lord, who
well understood astronomy, and who handled the subject with the learning
of a philosopher, was heard with impatience. With this example of the
advantage of being ignorant of the subject upon which he was about to
write, which he betters, by turning his real or assumed ignorance to the
account described:---he proceeds to give a history of alleged misdeeds
of workmen when combined---of many of their strikes, in which they are
described to be everything that is unreasonable, unjust, tyrannical, and
absurd, in the hope, as already in part quoted, that his readers will
believe of Trades' Unions that---

\begin{quote}
    \hspace{2em} Their aim and object is, in every case which we have
    been enabled to investigate, to \textit{stint} the action of
    superior physical strength, moral industry, or intelligent skill; to
    depress the best workman in order to protect the inferior workman
    from competition; to create barriers which no Society-man can
    surmount, and which few non- Society men dare to assail; and, in
    short, to apply all the fallacies of the protective system to
    labour. Such a system injures, first, the individual, whom it robs
    of a free market for his labour; secondly, the class of manufactures
    to which he belongs, by increasing the cost and diminishing the
    efficiency of the workmen; and, lastly, the nation at large, by
    curtailing the productive power, and consequently the wealth of the
    community.---P. 529.
\end{quote}

The only result of his ``investigation,'' as exhibited by himself, has
been the total exclusion of every fact and circumstance that would
justify, excuse, or extenuate the proceedings of the men, which, indeed,
appears to have been the only object he had in view in setting himself
to investigate the matter.

We have already shown the fallacies of the assertions, upon which this
confused argument rests, itself consisting principally of mere
assertion. The only part we have not alluded to that might be construed
into an application of the ``fallacies of the protective system to
labour,'' is the insisting, as Trades' Unions do, that all their members
should serve an apprenticeship, generally of seven years, to their
trade, and of sometimes endeavouring to limit the number of apprentices.
This last, however, is not now the general practice, but the exception.
It is always a needy or grasping employer who takes an inordinate number
of apprentices, and he never cares a straw for ``the nation at large,''
or ``the wealth of the community;'' he only aims, in thus using ``its
productive power,'' to enrich himself, no matter at whose expense.

The Act of Elizabeth, requiring an apprenticeship of seven years before
exercising a trade, has been repealed nearly fifty years, after having
been in force for two centuries and a half The ``Quarterly Review,'' in
its late article against Trades' Unions, states that this Act
``unquestionably exercised an important influence on English industry.
It stigmatised and punished the idle and the vagabond, directed the mass
of the people to manual occupation as affording the best means of
independent subsistence; and being acted on steadily from generation to
generation, it gradually educated a nation of skilled artizans.''---P.
480. And although the Act has been repealed so long, the practice has
not fallen into disuse. In some trades there is, probably, less to be
learned than in others; but there are very few trades in which long
experience and practice are not necessary, before the workman is able to
compete with those who are considered valuable to the employer. This
ability must be acquired at some period, and what more proper than that
in which the individual passes from youth to manhood---a time in which,
of all others, a facility of hand, and dexterity in performing difficult
processes, is the easiest acquired. This, then, being the most proper
period, and at the same time this also embracing the period when the
passions develope themselves, and are the most ungovernable, does it not
follow irresistibly that, if ever in a man's life he need the
experienced direction and restraint of his elders, it is this period? An
apprenticeship of seven years, from the age of fourteen to twenty-one,
whether considered necessary to learning a trade, as the best period for
acquiring it, or as conducing to the future welfare of the apprentice
himself, points itself out therefore as absolutely necessary. And this
is the reason why the repeal of the Act has made so little difference in
the practice, and why Trades' Unions, generally speaking, insist that
all their members should have served ``their time as apprentices.'' The
members are most of them fathers, and they know by experience the
beneficial effects of the practice, and the mischief that attends the
letting boys ``do as they like'' at this critical period. But we shall
be told this is no reason why the ``productive power, and consequently
the wealth of the community,'' should be ``curtailed.'' True. And the
same reason exists why the ``master villany of the earth, which includes
every other villany,'' and deadly sin---slavery, should perpetually
continue; for how else can ``the wealth of the community'' be so well
kept up by supplies of cheap cotton and sugar?

By the ``nation at large'' we suppose is meant what is termed the
``public.'' The public cares very little how things are produced, so it
gets them cheap, and the cheaper they become the more the public likes
it. The public, however, always respects those who so take care of
themselves as to prevent their suffering wrong. All the reasoning in the
world, nor all the admonitions of virtue, benevolence, or religion,
would prevent the public from preferring slave-grown sugar or
cotton---all the horrors of the slave trade notwithstanding---if it were
cheaper and as good as either grown by free labour. But, in justice to
the public it must be admitted, that it is also true, that, in the event
of men refusing to be slaves---supposing such a thing were to
happen---it would receive with a shout of mingled execration and
derision the wailing of those who complained that commodities were dear
because men were no longer enslaved to make them cheap. The reviewer,
therefore, has no occasion to make himself uneasy about the ``nation at
large;'' it never scruples to take care of itself, or fails to respect
those who do the same.

In giving his history of the alleged misdeeds of Trades' Unions, he has
not bettered his example. Nothing was ever overdone by Lord
Chesterfield. Indeed, on second thoughts we begin to doubt whether he
ever read the works of that nobleman. He proves too much by his
exaggeration; that is, if a jumble of alleged facts, disjointed from all
proof that they of necessity belong to Trades' Unions, can be said to
prove anything. And it is superfluous to say that he who proves too much
proves nothing. He wishes it to be believed that these misdeeds were
done by Trades' Unionists in and by their Union. If even they were all
true, to make them prove anything against Trade Societies it is
indispensable that he should also show of each and all that Trades'
Unions could not exist without these or similar enormities. This he does
not attempt to do; indeed he has entirely shut himself out from doing
so, by his real or assumed ignorance of their action---by his own
representation that they are secret, and that therefore nothing can be
known about them. He seems to have forgotten that his real or pretended
ignorance of their nature and action, so convenient for dark
insinuation, makes him incompetent even to argue the subject, much less
to prove anything, and consequently worthless as an authority. We have
already shown the utter fallacy of his insinuation, that Trade Societies
are secret, irresponsible in their government, or tyrannical; but that,
on the contrary, they are instituted for the purpose of correcting the
disadvantage inherent to the position of the working man when he
bargains \textit{singly} for the sale of his labour, namely, of being
compelled by his immediate necessities to take a price for it below its
fair exchangeable value. To assert that they, in any given instances,
have been turned to a bad purpose, is only to say of them what may with
truth be said of everything human.

The article consists of thirty-eight pages, out of which twenty-eight
are devoted to this history. Some fourteen or sixteen cases are
mentioned, which we suppose are by him deemed sufficient proof; and if
the careful suppression of everything alleged by the men as the reason
of their proceedings, with the aggravation of all that was alleged
against them, and the as careful setting forth as patterns of justness
and liberality the proceedings of the employers, could achieve a
demonstration of truth, the reviewer would be triumphant. It is rare to
find so much alacrity, so much \textit{con amore} readiness to accuse.
If the reviewer had been formerly an Unionist---a compositor, for
instance---and, like Hugh Miller, had been affronted by his former
associates, he could not be more bitter against Trades' Unions. His zeal
against them is as great and as indiscreet as that of an apostate. The
most absurd and exorbitant demands, upon the silliest pretences, the
most unjust and tyrannical proceedings towards each other, without sense
or reason, he charges, in the various instances he has given, against
the members of Trades' Unions. He refutes himself from the sheer
impossibility of such allegations to be true; or, which is the same
thing, to contain the whole truth. Our argument does not rest upon the
refutation of particular cases; it is, therefore, not necessary to
follow him through his twenty-eight pages of one-sided statement. It is
only necessary to take the first-mentioned, which is the strike of the
tin-plate workers, at Mr. Perry's, Wolverhampton, by which may be judged
how far he is to be relied upon for the rest. He revels in bitter
accusation with the vehemence of female malignity, incapable, like a
true virago, of seeing wTong in anything or anybody but the object of
her vituperation, recounting each accusation without pause or break in
the energy of her volubility, until her vehemence has fairly spent
itself. Such is the manner in which the subject is treated, and, as
might be expected, the conclusion is destitute of logical sequence.

After a page or two of uncertain maundering on the Hand-loom Weavers'
Commission, which he quotes, and with which he does, and does not, agree
in respect to alterations in the law of combination suggested therein,
which he very truly suggests ``do not amount to much,'' the conclusion
he comes to is, that there should be a parliamentary inquiry into the
state of industrial society in England, before a select committee. What,
an ``inquiry'' only after so much violent accusation about the
``mysterious tyranny,'' which has ``honeycombed'' English
society---eaten, as it were, into its very vitals---like that of
Continental nations, with its ``spies,'' ``enforced dumbness,'' "
\textit{suveillance}" and ``distrust on the right hand and on the
left?'' Why, he abjures himself. It is a confession that he has written
to the prejudices of his readers, and that he is frightened lest they
should seriously believe what he himself does not know to be true.
Though this is quite apparent from his mode of treating the subject, it
is singular that he should virtually confess to the fact. Yet with
strange perversity, like a virago, who, no matter how she has destroyed
her own credibility, ever returns to her first word, he reiterates as
his last word, this time by an assertion, in answer to a roundabout
question put by himself, that Trades' Unions are ``a secret organisation
of trades,'' which ``have undermined the groundwork of society in
England!''

We now proceed to show how far his version of the first of the cases he
says he ``investigated'' is to be depended on, by giving an account of
the same case, chiefly from the same authorities, which he himself
consulted in his investigation. The minutes quoted of the Committee of
the National Association, to which he had not access, only showing the
success with which the defendants were imposed upon, of the imposition
itself there was ample evidence in the report of the trial, which the
reviewer states he consulted.

Four pages are devoted to an account of the strike and subsequent trial
for conspiracy at Wolverhampton. The strike was at Mr. E. Perry's,
tin-plate worker, japanner, \&c., of that place, who was afterwards the
prosecutor.

Mr. Perry is said to ``have believed'' that he was paying ``ten per
cent, higher wages than the average of his trade in Wolverhampton; and
there seemed to be no conceivable reason why anybody should be troubled
about the affairs of the manufactory.'' According to the men, however,
Mr. Perry had been for years attacking their wages; and, at the time of
the strike, so far from his paying 10 per cent, more, he was paying 15
per cent, less than the average wages of the other manufacturers of the
place.

Mr. Perry admitted, in his evidence at the trial, that there had been a
strike in his factory in 1842; and, after some hesitation, that he had
agreed to a ``particular'' book of prices for his own shop, and that
``he had been distinguished in Wolverhampton by having a great deal of
opposition brought against his manufactory, by other persons not
connected with it.'' This does not much differ from the men's account of
the ``opposition'' to which he had been subjected through paying lower
wages than the other manufacturers of the town. With certain employers,
it is never their own men that are dissatisfied; it is always other
people who set them on. To rectify this, the men joined the ``National
Association for the Protection of Industry,'' and requested its
Committee to wait upon Mr. Perry in their behalf. Accordingly, Messrs.
Peel and Green, and afterwards Messrs. Green and Winters, in company
with ``three or four of his own men,'' waited upon him with a list of
prices, to which they wished him to agree. One of this deputation
afterwards became his foreman, and was one of the witnesses against the
defendants. The other employers, six firms in all, were also waited
upon, with the same list of prices, two of whom---Mr. W. E. Walton and
Messrs. Shoolbred, who were large manufacturers, and who employed more
men than all the others put together---returned a favourable answer;
one, Mr. Thurston, a conditional acceptance; the other three, Mr. E.
Perry, Mr. R. Perry, and Mr. Fearncomb, entered into a negotiation with
the deputation, or professed to do so. In reference thereto, Mr. E.
Perry, in a letter to the \textit{Daily News}, of August 20, 1851,
writes:---``I gave the parties to understand that I required time to
examine their proposal, and more especially to look into the book of
prices.''

In furtherance of this expressed intention, and to keep up this
appearance, Mr. E. Perry on several occasions met the defendants, and
actually agreed to the price of certain articles; at which time there
appeared to be the most perfect accord as to the negotiation. Mr. Peel,
as admitted by Mr. E. Perry, said that ``he was not for strikes, but for
peace,'' to which Mr. E. Perry replied, that ``if a satisfactory
arrangement took place, it would be through his (Mr. Peel's) good
management.'' Mr E. Perry also drew up a preliminary agreement
declaratory of such arrangement, which they that is, Mr. E. Perry, Mr.
Peel, and the rest of the defendants who were present, except
Rowlands---all signed.

How far the members of the Committee of the National Association were
justified in believing themselves, besides being appointed by the
journeymen tin-plate workmen of Wolverhampton, to represent them, also
fully recognised in that capacity by Mr. E. Perry, will appear by the
following extracts from the minute-book, written at the time, of the
National Association:

\begin{quote}
    \hspace{2em} ``April 9, 1850. At an interview that Messrs. Peel and
    Green had with Mr. Perry, it was agreed, in consequence of the
    meeting of the ironmasters then about to take place, to postpone all
    further proceedings for a fortnight; but still it was deemed
    necessary that the other employers should be personally visited, 
    Mr. E. Perry, in the meantime, promising to use his influence to obtain
    a conference with the other employers on the subject.''

    \hspace{2em} ``May 3, 1850. Mr. Green reported that Mr. E. Perry had
    consented to a conference of the employers and the men, one from
    each shop, together with Messrs. Peel and Green, as their advisers
    and advocates, to arrange the book of prices.''

    \hspace{2em} ``May 30, 1850. Mr. Green reported that Mr. Peel and
    himself had had two meetings with Mr. E. Perry about the
    preliminaries, and that the conference of the men and the employers
    would take place on Monday, June 3, 1850.''

    \hspace{2em} ``June 5, 1850. Mr. Green reported that the proposed
    conference between the employers and the men, commenced on Monday;
    that there were four employers and six men (one from each shop)
    present, in addition to himself and Mr. Peel, to arrange the book of
    prices.''
\end{quote}


It was at this meeting that Mr. E. Perry drew up the preliminary
agreement declaratory of such arrangement referred to above.

We give these extracts, not for the purpose of comment, but simply to
show that, whatever might be Mr. E. Perry's version of the affair, the
two members of the Committee of the National Association had good reason
to believe that he was perfectly willing to arrange respecting the
prices of the men. We may here mention, that, in the end, the two
employers, Messrs. Shoolbred and Co., and Messrs. W. E. Walton and Co.,
whom Mr. E. Perry, in his letter of August 20, calls the two ``youngest
firms,'' who, as before observed, employed more men than all the other
firms of the place put together, agreed to the terms of the men.

To return, however, to Mr. E. Perry. After keeping up this pretence of
negotiation, with all this apparent good faith and appearance of
agreement, for about three months, having employed this time in
``forwarding a large amount of orders then unexecuted; in accumulating
as large a stock of goods as possible, in anticipation of further
orders;'' and fortifying himself by what he called ``protecting his
men,'' in entering into written contracts with sixty of them, for one,
two, and three years, he threw off the mask, and ``frankly,'' as he said
in the same letter---``coolly'' we should have thought the most proper
word---told those with whom he had kept up this farce, that he never
intended from the first to agree to the list; that he was now provided
with men under contract, and should decline all further communication
with them.

It was under these circumstances that the strike took place, and no one
can wonder that it was carried on with feelings of bitterness on the
part of the men, especially when it was found out by them, so bound,
that the contract, while it bound them to Mr. Perry for its full period,
whether of one, two, or three years, and six months' notice afterwards,
was not binding upon him, who, by its terms, which they had never
strictly scanned, was able to turn them off at one month's notice. This
was the true reason why several of these men left him. Besides, they
considered themselves deceived. They imagined they would be employed
according to the list then in progress, while it was never intended to
be agreed to. This was the reason why the men, in the language ot the
reviewer, ``began to disappear;'' this was the reason why there were
``shop meetings.'' The men were irritated by the flagrant deception that
had been practised, both with regard to the prices and what they
conceived to be the injustice of their contracts. It is singular that
Mr. E. Perry should call this ``protecting'' his men. These men
ignorantly thought that contracts so inequitable, and consequently, as
they believed, so unjust, could not be legal. This foolish mistake was
soon corrected by the magistrate, who sentenced some to imprisonment for
breaking them. This, however, did not convince them that they were
morally binding, and they still continued to leave. To put a stop to
this practice, which was only bringing mischief upon themselves, the
``London meddlers,'' as the reviewer calls them, posted placards, urging
these men to return to their work; and in consequence many of them did
so.

The reviewer omits all mention or reference to the fact above mentioned,
that the two firms of Messrs. Shoolbred and Co., and the Messrs. Walton
and Co., agreed to the men's list of prices; and also, to the fact that
Mr. E. Perry, in his evidence on the trial of the men whom he indicted
for conspiracy, swore that he was in wages paying, to the best of his
belief, the same as these firms, except, it might be, that in £100 he
paid under 10s.; that is not quite $\frac{1}{2}$ per cent. less than
they did. We have already stated how much the men said he was paying
less. Taking his own sworn testimony, however, what motive could there
have been for the deceit he practised? This poor $\frac{1}{2}$ per cent.
would not, on part of the men, have stood in the way of an arrangement.
Our knowledge of Messrs. Peel, Winters, and green, the three ``London
meddlers,'' enables us to speak positively as to this fact. They would
only have been too happy to come to such terms. It must, if he is to be
believed, have been from the pure love of contention and strife.
According to his own words---and they are his words on oath---he had no
one to blame but himself for the strike which took place. As to the
charge against the defendant Green, of advising the apprentices to spoil
their work, though urged with great energy by Mr. E. Perry's counsel in
aggravation, it is perfectly clear the same as the rest. Mr. Green
denied it in an affidavit, and the Court believed him in preference to
the witness.

That there might have been some acts which came within the scope of the
law, considering the exasperation which the deceit practised upon the
men would naturally occasion, was to be expected; but there was no act
of violence during the whole affair. Nor was the only threat of personal
violence deposed to believed by the jury, as they on the trial
expressly state. There was some ``hard swearing'' on these trials---for
there were two: the first of the defendants who belonged to
Wolverhampton, and the second of the Committee of the National
Association. We find, by the \textit{Wolverhampton Herald}, of October
22, 1851, that Mr. E. Perry was brought before the magistrates of
Wolverhampton, to answer a charge of perjury oon the second trial,
brought by Mr. Peel, for swearing that he was seen by him (Mr. E. Perry)
in Wolverhampton on a certain day, and thereby, as there was no other
evidence against him, causing his conviction; Mr. Peel being then, and
nearly a month both before and after that date, undoubtedly in London.

As none of the defendants had the remotest idea that such evidence
would have been given, they were not prepared to rebut it, and it was
then impossible to bring evidence from London in time. The effect of
this evidence was to connect mr. Peel with a certain placard. The
present writer was shown at the time the ``Letter'' and other books
kept by Mr. Peel, in London, day by day. From the entries therein, it
appeared indisputable that the defendant was in London on, before and
after that date; he also has conversed with persons, who assured him
they saw and spoke to Mr. Peel, in London, on that day.

As, however, Mr. Perry brought two witnesses to swear that tey also saw
Mr. Peel there on that date, the magistrates, according to the report,
declined to commit. This decision, however, they said, would not debar
Mr. Roberts (the solicitor for the prosecution) from proceeding by
indictment. Mr. Peel, however, did not so proceed, for he very
reasonably said, ``I know that I was in London at the time Mr. E. Perry
swore I was in Wolverhampton; and so do others too; but how am I to know
that, when the indictment is tried, there will not be two more witnesses
to swear that they also saw me there on that day?''

The reviewer is very copious in his account of the alleged spiriting
away of the men engaged by Mr. E. Perry, and of the ``low intrigue,
deceit, drunkenness, and coercion,'' stated by the witnesses, some of
whom would drink to any extent while drink was to be had,---``Bumper''
Griffiths, one of them, who obtained this cognomen from his frequent
inebriety, being a specimen of this class. But the reviewer is entirely
silent as to the conduct of Mr. Perry, and of the likely consequence of
that conduct, unblushingly not only avowed, but made a merit of, and
triumphantly boasted of by him, when he threw off the mask. Is intrigue
only low, and deceit reprehensible, when alleged against workmen? The
fact was, whatever else there may have been, there was no deceit or
intrigue, properly so called, on the part of the defendants at all. The
men who did not go to work for Mr. E. Perry, or, having been engaged,
did not continue with him, perfectly well knew what they were doing, and
were willing to do it. Coercion was, therefore, out of the question. Mr.
E. Perry, by his own act, provoked the strike; and as, according to his
own sworn testimony, whatever might have been the real fact, he had
little or no advantage to gain by it, it must have been for the sake of
the strike itself. What else but a strike could have been expected,
after what had passed? If there is one thing more than another that
causes men to feel bitter towards another, it is a sense that they have
been deceived. What, therefore, could have been expected but that it
should have been carried on with bitterness---with not only determined
opposition, but with the contemptuous loathing that men always feel
towards him who they believe has deceived them? It may reasonably be
asked, if there were so little difference between the prices (according
to his sworn testimony) paid by Mr. E. Perry, and those paid by the
other employers, why did the men continue on strike? The men asserted,
so far from this being true, that he was paying about 15 per cent, less
than they did. And when it is considered how nearly flagrantly deceiving
his men is allied, from the very fact, to the intention of paying them
as little as possible, the statement of the men is the most likely to be
correct; in which case, it will be seen, that this gentleman, after
commencing his operations in this affair by the before-mentioned
deception, of which he triumphantly boasted, and carrying them on by a
series of inequitable contracts with his men, finished victoriously in a
court of law, by stating what opponents declared was incorrect. It is
not difficult to account for the apparent sympathy in this case of the
reviewer; for was he not himself occupied in suppressing half the truth
throughout twenty eight pages of good octavo print?

\section{THE ``QUARTERLY REVIEW.''}

Capital is said by this reviewer to fly turbulence and strife, and to be
timid, which is very true; but this is very incompletely stating the
question. Capital eschews no profit, or very small profit, just as
Nature was formerly said to abhor a vacuum. With adequate profit,
capital is very bold. A certain 10 per cent, will ensure its employment
anywhere; 20 per cent, certain will produce eagerness; 50 per cent,
positive audacity; 100 per cent, will make it ready to trample on all
human laws; 300 per cent., and there is not a crime at which it will
scruple nor a risk it will not run, even to the chance of its owner
being hanged. If turbulence and strife will bring a profit, it will
freely encourage both. Smuggling and the slave trade has amply proved
all that is here stated; and, what is worse, the slave trade and
slavery, by the profit attending them, have so debauched the public mind
where they exist---and be it remembered neither would exist but for the
profitable return of capital so employed---that it has turned the
so-called freest country in the world into a vast slave-pen; and, worse
still, it has set all the pulpits in the ``South'' of that country to
prove from the Word of God that this master crime is sanctioned by the
Almighty. This is what capital does, and is doing; which, alas for human
nature! no one can with truth contradict. Now, we do not mention this to
cast a stigma on capitalists; for we believe that the evil instinct
mentioned at the commencement, from which it all springs, is equally to
be found among our own class as the class that employs capital. Our
object is to show that, when labour and labourers are assailed by all
manner of folly, mischief, absurdity, and wickedness, being imputed to
them as being done by means of their Trades' Unions, that the same
things, and infinitely worse, are justly to be imputed to capitalists in
the employment of their capital; and that it is equally absurd to impute
to Trades* Unions any of these alleged vices as it is to impute to
capital all the crimes we have here enumerated. Capital in itself is
good, and so in themselves are Trades' Unions. They are a barrier
against undue reduction of wages, which is always prompted to increase
immediate profits where there is an immediate opportunity. Without them
that opportunity would frequently occur, of which advantage would not
fail to be taken. Both may be misused. Capital, we know, is, and our
opponents say Trades' Unions are; but that is no argument against
either. To charge the enormities above recounted against capital would
be as absurd as to charge the ``shillala'' with which a man's brains
have been dashed out, or the blunderbuss with which his head has been
riddled with bullets, with the murders committed by the use of them. The
same with Trades' Unions. We, therefore, except for mere illustration,
do not feel called upon to go into the different cases mentioned by
``Quarterly'' reviewer. Our argument does not rest upon their refutation
or denial. Concerning them, we only assert that they are one-sided, and,
consequently, do not convey the whole truth. That a Trades' Union may be
so worked as to be pernicious we no more attempt to deny than that
capital has also been employed for nefarious purposes. All we contend
for is, that such purposes are not the natural and proper uses of
either.

In the instances of strikes and other matters which the reviewer gives,
we see no reason to believe that he has, like his brother of the
``Edinburgh,'' imparted to them any bitterness of his own; he appears to
have taken them just as he found them, with all the abridgments fairly
made; but it must be always remembered that it is the employers'
versions that he has used. The men's version of the same transactions
would be very different. Take, for example, the Preston strike of 1854.
It is alleged that besides the 10 per cent, from the four employers who
refused the advance, which had been given by the other thirty-one or
thirty-two employers of the place, that the men wanted ``to fix a
standard of wages throughout the country, and it was upon this point
that the strike virtually took place.'' In reply, it is only necessary
to state, that the men from the first were anxious to refer what they
wanted to arbitration; an offer which the employers invariably refused.
This fact---which is a complete answer to all charges of unreasonable or
exorbitant demand on the part of the men---the reviewer does not
mention. The real state of the case appears to have been this: the firms
who had given the advance of 10 per cent, saw in this strike an
opportunity, by the process of a ``lockout'' of retracting it. When,
therefore, the men struck these four factories in the belief that they
were serving the employers who had given the advance, these very
employers were agreeing with these four employers to defeat the men, and
get back the 10 per cent, by locking the men out from all their own
factories. This was no doubt considered by them a masterly stratagem,
which must ensure success. The workmen, however, perceived it in a
different light; they viewed it as a coup worthy of Mephistophiles. It
was indeed what they did not expect; they never supposed, simple men!
that those, in whose behalf they believed they were contending, would
turn in this manner against them. It was their ``Leipsic,'' and what
their feelings were likely to have been in these circumstances may be
judged from those of the French army in that battle, when their allies
went over to the enemy and poured destructive volleys into their ranks.
Hence the exasperated and prolonged nature of the struggle.

It is quite true that according to bargain principle---under which all
contests of this kind should be conducted---perhaps their true course in
these circumstances would have been to refuse to contest the matter
farther. But who could expect them to do so? Men of higher grade, and of
far better education, have, under circumstances of less exasperation,
pursued what turned but to be a hopeless contest to the end and, if the
expression may be allowed, even beyond it---with, though losers, the
approbation of the whole community. But then they were not ``misguided
operatives.'' It must be admitted that they showed good fight, and
also---even according to the bargain principle---that they had, from the
long time it was doubtful which side would win, a good chance of
success. Mr. Cowell, at the Bradford meeting, stated the cause of their
defeat to have been the depression which arose from the uncertainty that
preceded the Crimean war, which was then imminent.

With respect to the fines alleged to have been imposed upon employers by
certain Trades' Unions in Manchester (p. 505), we can only say that we
hope the reviewer has been misinformed. This is the first time we ever
heard of such a thing. If true, it only shows that the evil and folly of
our nature will show itself everywhere. The men who counselled this,
would, if employers, be as unreasonable and as insolent as they, if
true, are here described to be as workmen; perverting the true use of
their Union, to insult their employers, just because they foolishly
thought they had the power to do so. Such men where such there be are
the bane of every society into which they come. Impotent for good, but
powerful for mischief, they destroy or spoil everything they touch. But
this description of persons are to be found everywhere, in every grade
of society; and it would be as unjust to decry Trades' Unions because of
their insolence, as to decry British merchants generally because of the
impertinence of the four Liverpool so-called merchants who lately wrote
to the Emperor Napoleon, ``requesting to know his intention.''

But it is said that the trade of shipbuilding has been driven from
Dublin by Trades' Unions and Strikes.

This charge, made by Mr, O'Connell, is founded on the evidence of a Mr.
Fagan, timber merchant (from 1818 to the time he was speaking), and Mr.
Morton, late a master shipwright (from 1812 to about 1828), both
ofDublin, given in the Blue Book on Combinations, 2nd Report, of 1838.
It does not appear that Dublin was at any time distinguished among the
other ports of the United Kingdom for the number and size of the ships
built there, or as its natural capabilities deserved. But this in
passing. Their joint evidence is to the following effect:---1st.---That
in consequence of the combinations the men did not earn their wages, and
that, therefore, the charges for shipbuilding v.-as so high, that no one
who could help it would have their ships repaired there. Ships that
wanted repairs were, therefore, only ``cobbled up,'' so as to take them
to Liverpool or other ports to be thoroughly repaired. Of the
ridiculously high charges an instance is given. 2nd.---That, in
consequence of apprentices being taken to remedy this, the men struck,
waylaid, and murdered one of them who was serving his time to Mr.
Morton, which was the ``finish of the shipbuilding trade.''

It is necessary to state that trade combinations, as described by these
and nearly all the other ``employer'' witnesses examined before this
committee, were not Trades' Unions in the proper sense of the word, but
always associated with violence; in short, as so many conspiracies ready
at any time to concoct and to perpetrate ``assassination beatings,'' as
Mr. O'Connell called them. Apart from violence, there seems, on their
part, to have been no objection to combinations. This also was Mr.
O'Connell's idea of them. According to these witnesses, violence and
intimidation was a sort of ``institution'' in Dublin and other ports of
Ireland, used by all who lived by labour; by brewers' draymen, canal
boatmen, porters in the egg-market, and others, whether belonging to
trade combinations or not. It is proper to state that all the
``workmen'' witnesses most emphatically denied that violence, or
incitement thereto, was ever done by their trade combinations. 1.---On
the first, Mr. Morton said, that while the men received 27s. a week,
``half of them were not able to do work worth 2s. a day; indeed, to some
of them he would rather pay 2s. a day to stop away.'' (Questions 5952 to
5955.) The instance of high charges is of a ship of 90 tons, which put
in for repairs; of which, he says, ``the shipwright's account for labour
was three hundred and odd pounds---that is, upwards of £3 (not less than
£3 5s. 6d.) a ton for repairing the vessel; and vessels can be bought in
the north of England for 30s. (for labour) a ton.'' (C)5959.)

In reply to this, it is only necessary to mention London, Sunderland,
Liverpool, Glasgow, and Dumbarton, among other ports where shipbuilding
flourishes, and is extending every year; there are what is called, in
this evidence, ``strong Unions.'' Surely it will not be said that the
incapacity or idleness described exists in either of these places. And
if not, the charge is irrelevant to Trades' Unions. Nor could it be on
account of wages; for in the port of London at this very time (1825),
the wages of a shipwright was 6s. a day, which rate continues up to the
present date. In respect to the instance given, as a charge, it refutes
itself In all the ports in the North of England there are the same
``strong Unions.'' If similar vessels, which in Dublin costs 3 5s. 6d. a
ton, for \textit{repairs alone}, could be bought out-and-out in the
North of England for £1 10s. for labour where there are ``strong
Unions,'' the charge is refuted out of his own mouth. We do not believe
these statements, but such is the evidence. 2.---On the second, Mr.
Morton says:---After two previous strikes, in 1812 and 1814, there was,
in 1825, another, about apprentices. He took one named Andrew Marchant.
``This lad (about 21 years of age) was beaten to death by a lot of men
at about half-past eight in the morning. He died in two hours.'' He
continued in business three or four years afterwards, when, because
trade was bad, and ``the number of his apprentices was so extensive,''
he quitted the business. This is the account he gives of the ``break
up'' of the shipbuilding trade in Dublin. The circumstances that led to
which, and probably the murder itself, it will be seen, took place
before the repeal of the combination laws. Indeed the acts of violence
before their repeal were greater than afterwards. Of this, Mr. Fagan
says:---``He recollected Mr. Morton; he was completely broken by
attempting to put down combination.'' Of this murder, he
continued:---``This man was murdered, he thought, about ten o'clock in
the morning; he could not exactly recollect the date; but that was the
complete finish of the shipbuilding.'' (3913 to 3921.) Shipbuilding had
not, however, ``been extensive'' in his recollection: he recollected
seeing two or three vessels on the stocks. (3922.)

From this account, it appears that the trade was ``finished'' by a
murder in the open day. The murderers all escaped, Mr. Morton said, in
such a manner, that there was ``no such thing as getting at any of
them.'' (5933.) When it is considered that this crime was committed when
``the bell was rung for breakfast,'' and, consequently, when there must
have been many persons passing and repassing from the Gasworks and Canal
Docks (where it took place being close to both), it is, we submit, most
incorrect, to call it a \textit{Trade} outrage. It was more like an
agrarian murder; it was done in the open day, by the connivance of all
passers by, who silently agreed to suffer its being done, and to allow
the escape of the murderers.

Of it we can only say, that, if this be a true account, the wonder is
not that shipbuilding was annihilated, but that any trade or manufacture
remained where such things could take place. We no more acknowledge any
connection between such acts as these and Trades' Unions, than we do
between a table knife and its use as a murderous weapon. Such acts and
violence generally, belong equally to all classes, and exist in every
place where there is impunity for them, whether there be many to one, as
in this instance, or on equal terms as in a faction fight. There are no
Trades' Unions in the United States ; yet where are there so many acts
of violence, from ``tarring and feathering'' up to murder? Even at
Oxford, where one would think such things would not be known, they occur
annually or oftener, in which lives have been lost, and will always
occur, while the authorities permit them ; for the way in which they are
dealt with amounts to permission. Indeed, by the grand display of this
pugnacity, on the arrival of his Royal Highness the Prince of Wales as a
student, it would appear these rows were also considered an
``institution,'' of which they ought to be proud. But wherever they
occur, and of whatever kind, for the continuance of such acts the police
are responsible, and it is only by the criminal negligence of the
magistracy that they are perpetrated. This was the opinion of Mr. Fagan,
who said that his life was threatened, and that for two years he carried
pistols (3876) in the daily fear of a murderous assault. He says that
the murder of Hanlon in the open day was because ``no person had been
taken up for any outrage before.'' (3954.) And again, ``If the
magistrates years ago had been active, and punished the cases brought
before them, I do not think the violence would have gone to the extent
it did.'' (3998.)

From the whole we think, that if these witnesses spoke the truth---and
they are the sole evidence---and that shipbuilding was finished in this
way, that it is not in any way to be laid to the use of Trades' Unions,
but to their misuse, and it is no more against them than---to use a
peaceful analogy---the misuse of capital in corrupting the electors of
Gloucester and Wakefield is against capital. But, after all, we doubt
the fact ; the same things happened there in other trades which still
went on, though doubtless all would languish in such a state of things.
It is most painful to read this evidence ; it is one continued recital
of wrong and violence from one to the other. Employer against workman,
and workman against employer. In the shipbuilding trade no workman was
examined, or there would have been a very different side to the
question. But in all the other trades there were. As the above evidence
is a specimen of what is alleged against the workmen, it is only fair to
give what the men alleged against their employers. They are accused of
the meanness of paying their men in bad halfpence, who, after being
bandied about from one to another, did not get good money for them.
(6540 to 6550.) The following account is an answer to the question why
they so ill agreed. The firm of Henry Mullens and McMahon, builders, is
spoken of: Their practice was, when they entered into a contract, to
engage about double the number of men the job would require, who were
generally sent for from the most distant ports, and, after they had
worked a short time at whatever rate they were engaged, about one-half
of them were sitting idle for the want of materials or some other
pretext ; but he held on in the hope of being again set to work, until,
far from friends and home, they were starved into obedience, when an
offer would be made to set them to work at reduced wages, which if they
accepted, the same or lower wages would be offered to those that were
kept at work with them ; and thus they would oblige them to reduce each
other to the lowest terms by which they could support a miserable
existence. This plan they pursued at the Dublin Docks and Library, where
those outrages commenced, and it was a portion of their own men,
rendered desperate by disappointment and want, who attempted to wreak
their vengeance on those whom necessity obliged to be the immediate
cause of their sufferings. (7948.)

Mr. Black will find, p. 209 of this report, his case of the two
``helpless orphans,'' quoted by him from a speech of Mr. O'Connell,
completely contradicted by the two ``orphans'' themselves ; and that
their employer, Mr. Mackie, not only ``never offered to take them as
apprentices, or offered them any protection,'' but that they served
their time under the protection of the men who were thus vilified ; and
that the said Mr. Mackie was a most avaricious employer, who used the
``truck system much to his own advantage.'' It would have been better if
Mr. Black had read this evidence before he gave his lecture.

There was one homicide by an employer, Mr. Carolin, who killed an
Unionist. He was tried and acquitted ; Mr. O'Connell, who was counsel
for the prosecution, said, in reference to it, ``that every avenue to
justice was choked.'' For the murder of Hanlon four Unionists were
executed. This was, doubtless, legally correct ; but the discrepancy
remained to rankle in the minds of those already exasperated. There were
strikes against what was considered to be an inordinate number of
apprentices ; they were used as a means of reducing wages, and resisted
on that ground. Many employers of all kinds agreed very well with their
men ; but there appears always to have been a certain class who were
never at rest unless they were reducing wages. Throughout the whole
evidence, amid the confusion of incidents and the wrong detailed on
either side---and on both sides there was much wrong---strange to say,
though employers deposed to ``assassination beatings'' and to going in
fear of their lives and to carrying pistols, there does not seem to have
been any difference in the course pursued. Perhaps they got used to it.
Each class appears to accuse the other of inordinate greediness : the
men of acting like the boy in the fable, who killed the golden goose ;
and the employers, of being ready from greediness to commit any act of
extortionate oppression. There were honourable exceptions, but the
leaven of evil was too strong to allow trade to flourish. It is
gratifying to add that this state of things has now passed away.

This reviewer has recorded the singular ``admission'' of Mr. O'Connell,
``that Trades' Unions had wrought more evil to Ireland than even
absenteeism and Saxon maladministration ;'' but he does not state when
this ``admission'' was made. Did it never occur to him that Mr.
O'Connell might, after hearing this evidence, have changed his opinion?
The reviewer could not be ignorant, that in the subsequent public life
of Mr. O'Connell not a single act was done by him that would justify the
belief that he retained it. But, on the contrary, if he did retain it,
he must have been the vilest of profligate political adventurers to
goad, as he did, to the verge of civil war, the passions of his
countrymen in the agitation for repeal which he set forth as that which
would be the salvation of his country, when he believed the master evil
under which she groaned was, in truth, to be laid to a very different
cause than her union with England. Whatever else may be said of Mr.
O'Connell, he was undoubtedly a shrewd man. He procured this committee
of inquiry, no doubt, in the belief of what he is said with singular
felicity to have admitted ; in the hope of putting a stop to the
frightful outrages---in fact, murders in the open day---which then took
place in Dublin and other places in Ireland, by not only artisans who
were Trades' Unionists, but by others whom we have mentioned, who appear
to have ``hung together'' without any specific ``union'' at all. The
cause of their continual repetition may be given in one word---impunity.
From that single cause we repeat---the Town and Gown riots at Oxford and
the outrages in the United States of America are also continued. And he
justly thought that if he could draw public attention to these acts,
which were a disgrace to the authorities of that city, boasting as it
did of a lord mayor with all the machinery of police authority, these
officers would be shamed by public opinion into exertions to put a stop
to them. In this he succeeded by this inquiry, for which intention,
indeed, more than one of the witnesses complimented him. But when he
heard the evidence on both sides, and how the violent acts were provoked
by equally violent extortion and oppression on the other, his opinion
must have been changed ; he must have seen that the evil was to be
attributed to a different cause, namely, not to Trades' Unions, but to
the provocation to which the users of these institutions were subject.
He could not fail to see that it was not the machine by which the thing
was done that was responsible for the mischief, but those who had used
it for this purpose. This view must have been entertained also by the
committee, for they made no report, nor appended a single opinion to the
evidence.

The cause of the building of ships not flourishing in Dublin is, we
think---if there be any value in the evidence here adduced---obvious
enough. Immediately it is, because they can be built and repaired
cheaper elsewhere. The cause of this cannot be in Trades' Unions,
because in the very places to which they go to be repaired, and in
which they are built, there are strong Trades' Unions. It cannot be in
the high wages that are paid in Dublin, because as high and higher wages
are paid elsewhere, where there are also strong Trades' Unions in which
the trade flourishes. It can scarcely be in the violence, for that has
occurred in trades that still went on. The case of the ship of 90 tons
costing £3 5s. 6d. per ton in labour for repairs, when a ship for
labour, from first to last, would only cost £1 10s. in the north of
England, if only partially correct, enigma as it is, is a perfect
solution why it is that no ships are likely to be built or repaired in
any place where such unaccountable charges exist. If the gentlemen who
have been so ready to denounce Trades' Unions were to read this
evidence, they would see that employers who act as this evidence
discloses they acted towards their men, and men who are provoked to act
as it also discloses towards their employers, are not on sufficiently
good terms, let them be in Ireland, England, or anywhere else. Trades'
Union or no Trades' Union, to ensure the success of manufactures or
commerce. No matter that labour is said by Mr. Buxton to be ``plentiful
and cheap,'' ``that there is abundance of water-power and first-rate
harbours.'' Capital and labour must go hand in hand ; each must be
willing to assist the other, or all is of no avail.

The predatory instinct of which we have spoken, and which belongs
equally to each, must be laid aside or kept down, that each may act in
harmony. If employers and employed use this instinct, ``as not abusing
it,'' the principle of exchange will flourish ; the employer will give
labour its just equivalent, and the employed will return it by good
measure pressed down, and running over. By no other course, where labour
is free, can manufactures and commerce flourish. A state of mutual
irritation, contempt on one side and hatred on the other, is
inconsistent with the state of mind essential to the progress of
industrial and manufacturing skill. The ability may be present, but it
will not be used to any successful purpose ; on the contrary, it will
either dwindle away, or take its flight to places more congenial to its
exercise.

\subsection{MASONS' STRIKE.}

As the ``Quarterly'' reviewer has recited many instances of
strikes to show the alleged animus of Trades' Unions, we quote the
following---given, it is hoped, with more fairness than his
instances---to show the alleged animus of certain of the employers. It
is abridged from the masons' ``Narrative'' of the strike, published in
1842.

Sir Morton Peto is reported to have said that he had, at the expense of
£16,000, defeated a strike ; and that he was ready to incur a similar
expense to support the ``declaration.'' This gentleman, who is equally
distinguished for his patriotism, his piety, and his success in making
money, has, no doubt, his own good reasons for being ready to repeat so
considerable an outlay. We make no complaint against him for having
pursued a course which he believed to be for his own advantage ; we only
claim an equal right for the men to do the same. The strike to which he
refers took place at the building of the Houses of Parliament, in
September, 1841. It was of the masons, against the tyranny of their
foreman, whose acts of brutality and profane language were made the
subject of complaint to his employers, Messrs. Grissell and Peto, July
13, 1841. There was also a complaint that there was not a proper supply
of certain tools, and that the men were forced to purchase their beer at
a particular public-house (the Horse-shoe and Magpie, Bridge-street,
Westminster). Upon these matters a deputation from the men had an
interview with Mr. Grissell on the ensuing Monday, when the following
complaint, as the most prominent, was entered into : One of the men had
received a letter, informing him of the death of his mother, from
Manchester, where she was then lying dead ; and, wishing ``to bury her,
dispose of her effects, and discharge her little liabilities,'' he
showed this letter to the foreman, and asked leave of absence for a week
or a fortnight, who informed him, with an oath, in reference to the
fortnight, that if he wanted his money---i.e., his discharge---he might
have it, but he would be d---d if he would keep a place open for
anybody for that time. The man then asked for a week, which the other
said was too long a time. He would be ``d---d,'' he said, ``if one day
to go to Manchester, one to bury her, and one to come back, was not
quite sufficient ; and he would be d---d if he would give any longer
time.''

Without obtaining leave the man went to Manchester, and returned at the
end of four working days to his work. Owing to the remonstrance of the
men he was not discharged. The foreman, however, would not allow him to
return to his usual work, which was of superior description, but put him
to the ``dirtiest, the coarsest, and the most inferior sort,'' which, it
need not be said, to a superior workman was the most effective mode of
insult that could be pursued. Mr. Grissell, after hearing evidence in
support of this complaint, which the foreman, who was present, did not
attempt to refute, affected not to believe it. All that could be got out
of him was, that ``the absurdity of the charge refuted itself.''

With respect to the tools, Mr. Grissell undertook partially to remedy
the matters complained of. With this the men, who did not wish to push
matters to extremity, were obliged to be content. With respect to the
beer, on which subject the men entered into ``somewhat fully,'' Mr.
Grissell, it is related, stated, in reply, that ``he could not make up
his mind to allow the men to have any other beer until he had consulted
Mr. Peto, his partner, inasmuch as there existed an understanding
between them and the brewers of that beer'' (p. 14 of the
``Stonemasons' Narrative''). Who the brewer was with whom there was
this understanding is not stated.

Now, if this statement is to be relied on, it is impossible not to
admire the skill and astuteness with which the gentlemen of this firm
turned every point connected with their business to account. They
endeavoured to make their men find their own tiles (for the soft stone)
; they are alleged to employ, and sustain in office, a brutal foreman,
who, on the same principle that selects the cruellest man that can be
found for a relieving officer ``to lessen the rates'' for the poor, is
supposed to extract the greatest possible amount of labour out of their
men ; and, finally, it is said, they endeavoured to turn to account, as
far as beer was concerned, their daily wants. No wonder the firm was
ready to oppose, at all risk, any alteration in an arrangement so
beneficial to themselves. Perhaps, after all, £16,000, which at 5 per
cent. is only £800 a year, was well laid out in the endeavour to sustain
it. For, if by these appliances sixpence a day additional could have
been got out of each of the men, 230 in number, it would have amounted
to £32 10s. 0d. a-week, or £1,690 a-year ; if only threepence, it would
have amounted to more than the interest (£800) at 5 per cent. of this
sum.

Many will exclaim that in this view of the question Messrs. Grissell and
Peto did right in not discharging so profitable a servant, as it was by
his means that so large an amount was obtained that they were justified
in retaining him at any risk. Exactly so. But by the same rule, and for
the same reason, namely, that they were made to produce an extra
quantity of work without being paid for it, and therefore by the true
principle of exchange---to say nothing of his profane brutality---the
men were justified in refusing to work under him.

The men on the whole believed that this interview would have a good
effect, and resumed their work. For about six weeks there appeared a
change for the better. This, however, did not last, for the foreman
reverted to his old courses, with aggravated violence. With respect to
the beer also, with or without the direction of his employers, finding
that the men drank water instead of the beer in question, he locked up
the door leading to the pump. The following instances were given of his
cruelty and tyranny, to which the men refused any longer to submit :

One instance was of a labourer who fell from a scaffold in their employ
and broke his leg, who, as soon as he was able, returned to his work.
The foreman told him to go about his business, as he did not want lame
cripples like him. This case was controverted. In an investigation which
took place, however, though there were some mistakes as to the
individuals mentioned, the fact remained unshaken.

Another, that of a man who was taken ill, and obliged to leave his work
for a few days. When he came back, expecting to resume his employment,
the foreman, seeing him not look so well as previously, told him to
leave the ground, as he, the foreman, wanted none but sound men there.

Another was of a man who had obtained leave of the under-foreman to
remain at home, as his wife was dying, and who died the same day. When
the man came in the morning to resume his work, he was told by this
gentleman to ``go and die with her, and be damned'' (p. 25), and
discharged him accordingly.

They therefore determined no longer to work under such a ruffian, and
gave a written notice to Messrs. Grissell and Peto that they would leave
unless he was removed. This led to a correspondence with the
Commissioners of Woods and Forests, under whose superintendence the
Houses of Parliament were then building, both on part of the firm and on
part of the men, which ended in the said commissioners believing the
employers and disbelieving the men, and acting accordingly---which, we
regret to say, is too often the case in such matters.

When Messrs. Grissell and Peto received the notice from the men that
they would leave, they stated that they had ``every disposition'' to
remove any real grievance, when ``properly represented ;'' but they did
not see that they were warranted in removing their foreman. They also
regretted that the men did not ``sufficiently appreciate the
consideration they had always experienced from their hands.'' It
certainly does seem remarkably cool, after what had passed, that this
firm should complain that their men did not ``sufficiently'' appreciate
the ``consideration'' they had received. Did they think the men ought to
have remembered their kindness in providing them with beer on an
understanding of their own with the brewer? The men, 230 in number,
left on the 18th. A deputation, at the request of the firm, had an
interview with Mr. Grissell, who repeated with much earnestness the
promises made in their letter. Upon the men reminding him, however, that
similar promises had been made before which had not been fulfilled, he
changed his tone, and tearing the workmen's letter of complaint in
pieces before them, he most emphatically declared that to complain of
his foreman's conduct to him was perfectly useless ; that on works so
extensive as those as he (the foreman) superintended, it was necessary
that a discipline similar to that maintained in the army or navy should
be enforced, and the masons, in not quietly submitting to it, must be
bent on ruining the firm ; that the competency of the foreman to
superintend such work would naturally consist in his being such a man as
he was ; that a ``college-bred, educated man would not suit him ; a
harsh and severe individual would only answer his purpose'' (p. 31).

The relations of buyer and seller of labour, of master and servant, the
subordination of the sailor or private soldier to their respective
officers, and all but the slave to his driver, are here jumbled together
in this gentleman's mind ; but, amid this confusion of ideas, and,
probably, its cause, there was one chief thing, namely, the profit of
the firm, to which the whole tended, from which there was no swerving.
All remonstrance was therefore ``dictation'' as to the ``manner in which
they were to conduct their business.'' As this firm adopted the conduct
of their foreman, the masons employed on all the other works upon which
they were engaged left their employment, numbering in all to about 400.
The strike lasted from September 13th, 1841, to May 25th, 1842, at the
end of which period these 400 men, by obtaining employment elsewhere,
were reduced to 80, when, ``from the open support given by the
Government to the employers, in receiving from them inferior work and
materials, and extending the time originally allowed for completion,
little hopes remained of removing the foreman'' (p. 86). The remaining
number resolved ``to relieve their fellow-workmen of supporting them''
by quitting London in quest of employment elsewhere. During this period
(nine months), the ``Narrative'' states that, of the 400, only five
cases of ``traitorism'' occurred ; ``during which, also, not a single
breach of the law was committed, or a police case occurred.'' The masons
did not succeed in their object, it is true, but then they did without
Messrs. Grissell and Peto, and these gentlemen did without the united
masons, and so the strike ended.

\subsection{BUILDERS' STRIKE AND LOCK-OUT.}

Of the nine hours' movement, we are of opinion, though
desirous of shortening the hours of labour, that it is one of those
things which do not admit of being decided by the issue of a strike. The
resolution for the strike had been carried, not by the leaders, but by
the body at large, by ballot. The time and place, however, had not been
determined ; and consequently, in the absence of further irritation, it
might have been postponed indefinitely. It might never have taken place
at all. The master builders, however, hastened to supply the irritation
requisite, by discharging the four men who each presented a memorial on
the subject to four firms, who had been addressed separately. Still, the
men hesitated until the man had been discharged by Messrs. Trollope and
Son, when the strike took place. What followed is so well known that no
further notice is required here. The master builders replied by a
``lock-out'' from nearly all the firms in the building trade. This is
said to have been done on a wise calculation. The lock-out was continued
for a month unconditionally, by which it is clear that the basis of such
calculation was revenge. If there had been any proper calm calculation
they would have discovered that there was, as we have shown, no
necessity for the lock-out at all. They now seek to break up trade
combinations in their trades altogether. Among the many extravagantly
foolish things said by the ``Edinburgh'' reviewer, he compared the
position of the members of the Amalgamated Society of Engineers as being
``very like the case of the Maryland or Washington Negroes let out on
hire, or Russian serfs under the obrok system'' (p. 552). As the
applauded advocate of the ``declaration'' and of the lock-out speaks of
slavery and serfdom, it may be worth while to inquire in what they
consist.

We have spoken of the principle of exchange being the root of civilised
society ; that, in short, civilisation is only the putting into practice
the four elements of which it consists. Liberty cannot exist without the
right to the use of this principle. A slave cannot as a right demand
anything in exchange for his labour ; neither can a serf. The power of
free exchange in either would immediately alter their condition to that
of free men. It follows, therefore, that whatever reduces the power of
demanding a right and proper equivalent for his labour, or, indeed, for
anything else he may have for sale, brings that man in just such
proportion near to one or other of these conditions. Hence the abject
state of those whose necessities compel them to take anything offered for
what they have to dispose of. And if a man be compelled by another to
exchange or sell his labour to him at a disadvantage, he is deprived, in
so far, of his natural liberty ; and he is to all intents and purposes a
serf to him who so compels him.

Indeed, in the enforcing the ``declaration'' there is very little
disguise in the matter. The belonging to a combination makes the men
``too independent,'' ``insubordinate,'' ``mutinous ;'' all of which
expressions indicate a state of serfdom, not of free exchangers, and of
approval thereof. And so long as this compulsion is used, and in so far
as it extends, the man is to all intents and purposes a serf ; and he
only is not used so in other respects because such use is impossible in
a free country.

The reviewer admits slavery and serfdom to be as we have stated, only he
contends that, what he calls the control of the combination over the
sale of their labour reduces the men to this condition. But he forgets
that the men use the combination---which he so curiously distorts---to
give them the power of correcting a disadvantage, namely, their
immediate necessities when taken singly, that they might exercise this
right of free exchange with benefit to themselves ; while the
``declaration'' by destroying their combination would take away this
power, and consequently transfer this benefit from them to their
employers. The one is a voluntary arrangement for their own advantage,
the other a compulsory one, at their expense, for the advantage of their
employers. The reviewer here, besides confounding freedom with
compulsion, mistakes the gaining of an advantage for the losing of one.
But as the power of free exchange lies at the root of all liberty and
right, so the protection of it is the foundation of all law. There is,
therefore, an inconvenience in a free country, where this is
acknowledged, attending such enforcement. As no man has a right to
prevent another even indirectly, much less directly, from doing that
which is lawful and beneficial to himself, there are certain awkward
conclusions to which the enforcers lay themselves open. For example,
pursuing the same argument, it is said that the lock-out is a great
movement, at an immense sacrifice, in favour of ``free trade,''
whatever that may mean, in labour. Free trade is an incomplete
expression ; let us, therefore, adopt the more comprehensive term, free
exchange, which includes all honest trade. An exchange, to be free,
supposes both parties to be on somewhat equal terms ; otherwise, it is
likely to become swindling or robbery. Advantage taken by one of the
other's ignorance, weakness, or necessities, will very soon become one
or other of these crimes. For illustration, though one is scarcely
needed, suppose an exchange of blows. It would not be a free exchange if
one of the parties were to have one arm disabled or tied down while the
other had the free use of both. To take such advantage would be, and is
justly deemed, cowardly and infamous. By the ``declaration'' the
employer wishes to deal with his men singly, so that he, whenever he
pleases, may give the ``sweaters''' price for their labour ; their right
arm as bargainers being tied down by their necessities in its sale. This
he calls free trade, but the freedom is all on his own side. Call it
trade, if you will, it is not free exchange. This state of things he
wishes to perpetuate, that his workmen may never be able to prevent him
taking this advantage ; never be on anything like equal terms with him
in this transaction of exchange. He therefore wishes for his own
benefit---though in a different form---the result commonly following the
two peccadilloes mentioned above, which he is quite aware, also follow
an exchange that is not free.

On the cool insolence of refusing to allow their men to combine, when
they are in strong combination themselves ; of the declaring the men's
combination illegal, when their own is pronounced by high and hitherto
uncontradicted legal authority to be illegal, and other drolleries, it
is not necessary to comment.

\subsection{APPLICATION OF THE FUNDS OF TRADE SOCIETIES.}

The ``London Society of Compositors'' is a Trade Union ; 2,200
paying members ; re-established, January, 1848, ``to protect the wages
of labour, agreeably to the provisions contained in the London Scale of
Prices (agreed to by a conference of Master Printers and Compositors in
1847), the Scale of Prices regulating News and Parliamentary Work, and
such customs and usages as belong to the profession, not directly
mentioned in the said Scales.'' The terms of admission are, the serving
of seven years to the business, either by indenture or patrimony. The
weekly payments of each member are ``proportioned to his earnings---
i.e., under 10s., nil. ; 10s., and less than 20s., 2d. ; 20s., and
less than 30s., 3d. ; 30s. or more, 4d. per week.'' About 700 members
pay an additional 2d. per week to what is called the ``Provident Fund,''
which entitles them to receive a small weekly allowance when out of work
; and on an average of the last ten years, the Society has paid £500 per
annum to unemployed members. The Society has a library and reading rooms
at 3, Raquet-court, Fleet-street. The library contains more than 4,000
volumes, about 1,000 of which are in constant use. The reading-rooms are
supplied with all the principal newspapers and periodicals of the day,
and are open daily from 9 a.m. till 10 p.m. The librarian, within those
hours, will show any respectable person over the premises. The Society is
also the medium of considerable sums being raised for the benefit of
widows and children of members in distress ; and it subscribes yearly to
a dozen or more of the principal medical charities of the Metropolis,
that it may be enabled to give letters of recommendation to sick members
or their wives or children. A sum is allowed members whose tools may
have been destroyed by fire. The library and reading room were
established (in 1855) to prevent resort to public houses.

Our own Society---the Bookbinders', established 1787---is similarly
constituted, having also a library of about 1,800 volumes, which, for
some time, was the only one connected with a Trade Society, and an
office, a house of call apart from a public-house, at No. 5,
Raquet-court, Fleet-street. The terms of admission are the same. Number
of members about 600. Allowance to members out of employ 8s. 6d. a-week
for ten weeks ; and sometimes there is a subscription to give additional
relief to their unemployed members. Payment to the Society, 2s. a-month.
Relief to the sick is granted by a levy of 6d. a member upon a
``Petition'' for such benefit. There are also six pensioners, two at 3s.
a-week, and four at 2s. 6d. Amount paid to unemployed members, for the
last seven years, £3,055 2s. 6d. In disputes on wages, prices, \&c.,
during the same period, £98 14s. 4d. ; which will show, in ordinary
circumstances, the actual proportion that ``Strikes'' bear to the
proceedings of a Trades' Union. For the sick, by ``Petition,'' has been
paid, for the same period, £1,309 1s. To the library £10 a-year is
applied for the purchase of books.

These are two Trades' Unions pure and simple ; their benefits for sick
being either voluntary or by additional levy. The amount paid by the
compositors to their ``Widows and Orphan Fund'' is quite as much, or
more, in proportion to number, as that paid by our own Society to the
sick. We are not aware that any more is here put down than is paid by
any other Trades' Union. By these, as a sample of the rest, the
``Edinburgh'' reviewer may see how their ``Capital'' is disposed of.

\subsection{``DICTATION.''}

Trades' Unions are brought to bear against employers in
disputes for wages. They assert the principle of exchange. Now, without
insisting that in all cases they assert this principle wisely, this is
the inherent source of their action. They say such and such labour is
worth a certain price, which is either agreed to or not by the employers
who have need of that labour. Now, how is this demand oftener than not
met by opposing employers? By the reply that they will not be
``dictated to,'' or ``interfered with in the way in which they see fit
to conduct their business.'' It is plain, therefore, that into the minds
of those who use these expressions the idea of labour being subject to
the principle of exchange does not enter. With them it is something
which should always be at their command at their own price, resistance
to which is ``insubordination'' and ``mutiny.'' We should be glad to be
informed wherein the ideas of such employers differ from those of serf
owners in like circumstances? Would Mr. Black inform us what would be
the probable conduct of the sellers of corn or sugar if they were
accused of insubordination and mutiny, and treated accordingly by their
customers, because the said sellers would not sell at the others' own
price? And yet such, when reduced to its true issue, is the idea of the
``Edinburgh'' reviewer and the numerous class he represents. He and they
would put down Trades' Unions, because they, as he asserts, impede the
action of those who employ labour---they ``dictate to'' and ``interfere
with them in the way in which they see fit to conduct their business.''
Now, conceding to these, and to all employers, the full right to refuse
to agree to the wages their workpeople demand, we ask, is there anything
in the known acts of this class so pure, guileless, and free from blame
for oppression, cruelty, or dishonesty, that they should be exempted by
law from such ``interference,'' and, as they are pleased to call it,
``dictation?'' Were not certain proprietors of coal mines, who employed
females in a state of nudity to work with men in their mines, so
impressed with the idea that such was the proper mode of conducting
their business, that it required an Act of Parliament to convince them
of the contrary? Were not a majority of the cotton manufacturers so
thoroughly impressed with the idea that the employment of children of
tender years, and for excessive hours, was the only fit and proper way
of conducting their business, that it also required an Act of Parliament
to restrain them? Was not the truck system---that is, the payment of
wages to a great extent, sometimes wholly, in goods, very often of bad
quality at the price of good articles, sold by the employer---also
deemed a fit and proper mode of conducting their business by many
employers, until an Act of Parliament was required for its repression?
To give all the instances that might be adduced to show the contrary
would more than fill this publication ; take, therefore, only the most
recent instance in the Hinckley stocking weavers, the employers of
which trade let out the frames, which cost £3, at an annual rent of
£2 12s. (1s. a week). In the trial, Homer v. Taunton, wherein this
came out, it was, as reported in the \textit{Standard}, of Dec. 29,
1859, stated on oath, illustrative of the way in which this rent was
exacted : ``That, when trade was bad, witness only earned 1s. 10d. a
week, and then his employer took off the full frame rent of 1s., and
left him 10d.'' It is true that, after much remonstrance and difficulty,
he got off that week by paying only sixpence out of the 1s. 10d. ; but
even that shows how rigorously this rent is exacted.

Now we do not blame a class for the acts of a certain number of
dishonest persons who may happen to belong to it. Our object is to show
that the predatory instincts of our nature belong to the employer class
as well as to every other ; and that, therefore, they have no right to
claim exemption from the checks to which they may be subject from their
workmen in the matter of wages ; that there is no more reason to suppose
that their estimate of what ought to be paid will be just any more than
that of their workpeople.

It has been stated, in the \textit{Standard} of November 24, 1859, for
the information of the ``misguided operatives,'' that the employers have
to struggle against hardships and difficulties no less severe than those,
though of a different kind, endured by working men ; instancing in times
of panic, when wages are not reduced, that profits are swallowed up in
high discounts, when they work not for themselves, but for their bankers
; that they also suffer from bad debts, Mr. Commissioner Fane having
stated that upwards of £50,000,000 annually, or £1,000,000 a week was
sacrificed in bad debts ; and as the employed do not pass through the
Bankruptcy Court, the ``loss fell on the class of employers.'' The
writer of this article might also, with great propriety, have said, for
the same reason, that this loss was inflicted on the ``class of
employers,'' not by working men, but by one another. Talk of the expense
that now and then a ``strike'' involves, or the far greater expense of a
``lock-out,'' what is the expense of either, or both put together, to
this enormous loss inflicted upon the employer class by the recklessness
or dishonesty, or both, of themselves? It certainly is very cool in the
writer of this article to suggest, as he does, that their workmen ought
to make up, by a reduction of wages, this wholesale plunder of each
other. The banker will not forego his claim ; so, from the nature of
things, the workmen must, as there is nothing left out of which he can
be paid his proper wages. It is also very cool of this class, after
robbing each other to so enormous an annual amount, to assume a virtuous
indignation, and tell their workpeople that their demand of a certain
wage is a dishonest one, and that the combination by which such demand
is supported shall be put down, and to issue a declaration that they
will employ no one who belongs to it.

The \textit{Times}, of December 20, 1859, has the following in a leading
article :

\begin{quote}
Well, let us see, the British manufacturer has his virtues, as we know
; his home virtues we will say nothing about, lest Sir Cresswell
Cresswell arise to rebuke all universal eulogy. But he also must have
his code of trade virtues. He must not forge a trade mark ; he must not
send an article out of his warehouse which would kill or maim any one
who should use it ; he must not sell a wooden stick and guarantee it to
be cotton ; he must not sell cast iron cutlery and call it steel ; he
must not send out a pair of scissors never intended to cut, or an axe
that would fly to pieces at the first stroke. If there are people who
have in the trade a better name than himself, he must never outstep the
fair pace of emulation ; it must not enter into his mind to forge their
names and destroy their credit by affixing their brand to coarse and
worthless goods. These are the very rudiments of commercial morals.
They are the equivalents to ``Thou shalt not kill,'' ``Thou shalt not
steal,'' ``Thou shalt not bear false witness.'' Yet, heaven help us !
these are the very acts which are attributed to these very respectable
persons as every-day practices. These are cropping up in our law books
and taking rank in our police cases. They are growing into a ``custom
of manufacturers.'' They have been solemnly presented to a court of
justice for its sanction.
\end{quote}

We make no comment on this extract. We do not cite it, or mention
anything else belonging to capitalists as a \textit{tu quoque}, or as
a set-off against any misdeeds of which the working class may have been
guilty, still less to draw attention to their vices ; for it is written,
``There is none that doeth good, no not one.'' We only ask whether
capitalists have any right to claim exemption from the natural check
which the assertion of the principle of exchange by combinations gives
to their workpeople? or whether it would be conducive to the moral
prosperity of the community that the law should, as many of this class
desire, so exempt them?

\section{CONCLUSION.}

In the foregoing pages we have attempted to show that Trades'
Unions are necessary to secure to the workmen a free, and consequently
fair, exchange for his labour, which, from a disadvantage inherent in
his position, he is unable otherwise to obtain. We have also---we hope
successfully---replied to the various charges brought against them. In
doing so we have endeavoured to trace out the cause which leads to the
misuse of capital on the one hand, and of Trades' Unions on the other.
This we find to be an instinct of evil inherent in the very nature of
man, belonging not merely to capitalists and workmen, but to the whole
human family, which, for want of a better term, we have called the
predatory instinct. Within due bounds this instinct, doubtless, lies at
the root of all energy and enterprise ; allowed to predominate, it is
the source of every evil. To subdue this to its proper use should be the
aim of all who would promote the well-being of their own and every other
class of society.

Applying this to employer and employed, it is of the highest importance
that there should be a good understanding between them ; that neither
should vex or offend the other. Both are afflicted equally with this
evil instinct, and each should learn to bear with the other. Neither
class can injure the other without at the same time injuring itself.
Both are so essential to each other, and so intimately connected, that
one cannot be injured without the other feeling it ; and, consequently,
no triumph can be gained at the expense of justice by either class over
the other with impunity. While the workmen should weigh well the price
they put upon their labour, and endeavour to make themselves acquainted
with the principles that determine its value, the employers should act
reasonably, and listen to what their men have to say.

The ``Quarterly'' reviewer, whose candour we have before acknowledged,
well observes, p. 521 :

\begin{quote}
``At the same time employers ought not to stand too strongly upon their
rights, nor entrench themselves too exclusively within the circle of
their own order. Frankness and cordiality will win working men's hearts,
and a ready explanation will often remove misgivings and dissatisfaction.
Were there more trust and greater sympathy between classes, there would
be less disposition to turn out on the part of men, and a more
accommodating spirit on the part of the masters.''
\end{quote}

The italics are ours, as the endorsement of forty years' experience ;
for it is certain that the working classes will give their esteem and
confidence to those whom they believe to be actuated by similar feelings
towards themselves.

We have before said the true state of employer and employed is that of
amity, and that they are the truest friends, each of the other---for
each derives his revenue from the other. And the fact is,
notwithstanding all that has been written to the contrary in the
``Quarterly'' and ``Edinburgh'' Reviews, and from them down to the
humblest writer who has added his little spark to the general blaze---
that this state is for the most part their actual condition. Under no
other circumstances could the trade and manufactures of the country have
so greatly prospered and extended. It should, therefore, be the duty of
both to prevent this harmony being interrupted. Each should consider
this state their true relation, and consider its interruption the
greatest of calamities.
